\section{Introduction}

A post-quantum rollup architecture has two orthogonal concerns:

\begin{enumerate}[leftmargin=2em,nosep]
\item \textbf{The rollup-batch envelope pattern} -- how a sequencer
  composes transactions, threshold-signs a state-root tuple, and
  posts a single \texttt{RollupBatchTx} that the parent L1
  verifies inside normal block validation. This is the
  contribution of LP-218: the wire format, the
  parent-L1-verifies-the-sig pattern, the inherited cert mode,
  and the first P3Q kind ($0x01$ Pulsar).
\item \textbf{The per-kind verifier specifications} -- the
  individual cryptographic primitives that can occupy a P3Q
  \texttt{kind} slot, their wire layouts, their gas calibrations,
  their security assumptions, and the operator semantics for
  composing them into a cross-family-maximum PQ posture. This is
  the contribution of LP-220.
\end{enumerate}

LP-218 fixes the pattern and the slot. LP-220 fills in the
remaining two kinds.

\subsection{Why a single slot}

A naive design would allocate one EVM precompile slot per PQ
family: \texttt{0x012205} for Pulsar, \texttt{0x012206} for
Corona, \texttt{0x012207} for Magnetar. The earlier draft of
LP-220 did exactly that, treating the rollup-batch wrappers as
sibling primitives. Three slots, three ABIs, three audit
surfaces, three call-paths in EVM.

The 2026-06-03 decomplection collapses this into one slot with
kind-byte dispatch. The motivations are not aesthetic:

\begin{itemize}[leftmargin=2em,nosep]
\item \textbf{Solidity callability is a UX surface.} A rollup
  contract that wants cross-family-maximum verify writes three
  \texttt{P3Q.verify(kind, ...)} calls at a single address. Three
  slots would mean three address constants, three ABI imports,
  three precompile registrations to track in the operator's head.
  One slot keeps the developer surface coherent.
\item \textbf{Audit surface compounds.} Three slots with three
  wrappers means three independent code paths around the
  threshold-signature verifier core -- the wrappers carry their
  own structural parsing, length checks, gas billing, and
  domain-separation contexts. Each is an audit boundary. One
  slot dispatches into per-kind verifier cores via the
  \texttt{P3QVerifier} interface and leaves only the dispatch
  table to audit.
\item \textbf{Future kinds add a byte, not a slot.} If a fourth
  PQ family enters the Lux stack (a hypothetical code-based
  signature with HQC-style structure, say), it lands as one new
  \texttt{kind} byte in the dispatch table and one new
  \texttt{P3QVerifier} implementation. No new precompile slot
  allocation, no new Solidity function shape, no new audit
  boundary at the EVM-side.
\item \textbf{One and only one way per LP-200.} The ZAP Stack
  umbrella codifies the principle: every distinct concern has
  exactly one canonical entry point. ``Verify a rollup-batch
  threshold signature on-chain'' is one concern; the kind byte
  selects which primitive answers the question.
\end{itemize}

\subsection{The three kinds at a glance}

\begin{table}[h]
\centering
\small
\begin{tabular}{@{}llll@{}}
\toprule
\textbf{kind} & \textbf{Name} & \textbf{PQ family} & \textbf{Underlying scheme} \\
\midrule
$0x01$ & Pulsar   & Module-LWE   & FIPS-204 ML-DSA-65 threshold \\
$0x02$ & Corona   & Ring-LWE     & Corona two-round threshold \\
$0x03$ & Magnetar & Hash-based   & FIPS-205 SLH-DSA threshold \\
\bottomrule
\end{tabular}
\caption{The three production kinds at slot \texttt{0x012205}.}
\label{tab:kinds-overview}
\end{table}

A simultaneous structural break of all three would require
independent advances in (i) Module-LWE cryptanalysis,
(ii)~Ring-LWE / ring-structured lattice cryptanalysis, and
(iii)~collision resistance of the SHA-2 / SHA-3 family. These
three assumptions are widely understood to be mutually
independent. The cross-family-maximum posture LP-217 PQ-heavy
codifies relies on that independence; LP-220 supplies the
on-chain machinery to actually compose all three kinds inside a
single rollup contract.

\subsection{Honest measurement state}

Throughout this paper, performance numbers are tagged
\textbf{measured} or \textbf{projected}.

\begin{itemize}[leftmargin=2em,nosep]
\item \textbf{Measured} numbers come from the
  2026-06-04 Quasar 4-scheme matrix bench
  (\texttt{lux/threshold/protocols/parity/parity\_test.go}) on
  M1 Max (10-core darwin/arm64) and Spark (NVIDIA GB10 Grace
  Blackwell, 20-core linux/aarch64). Numbers are wall-clock per
  per-validator signing round at $N{=}\{16,32,64\}$, mean of 50
  rounds.
\item \textbf{Projected} numbers are CPU verify costs scaled by
  GPU speedup ratios from kernel-level micro-benchmarks. GPU
  dispatch is not exercised by any path in the matrix bench --
  \texttt{nvidia-smi dmon} mean \texttt{SM\%} during the 2-hour
  Spark run was \SI{11.9}{\percent} (idle/background only). Two
  boundary gaps prevent dispatch: \texttt{lux-lattice.pc} is not
  installed on Spark (\texttt{find\_package(MLX REQUIRED)}
  fails), and \texttt{lux\_slhdsa\_*} entry points are not yet
  exposed in \texttt{lux/accel/c\_api.h}. Until both close, all
  GPU latency claims in this paper are projections.
\end{itemize}

The architectural decisions (slot, ABI, dispatch, parameter set
choice) are not contingent on the GPU dispatch closing. They are
contingent only on the CPU verifier correctness, which is byte-
equal to the underlying NIST-standard (Pulsar/Magnetar) or to
the Corona reference implementation (Corona) and is wired
already.

\subsection{Roadmap}

Section~\ref{sec:architectural-decision} documents the
2026-06-03 decomplection in detail: the prior sibling-slots model,
the unification onto \texttt{0x012205}, and the relationship to
LP-0120's separate raw-primitive verifier slots.
Sections~\ref{sec:pulsar-kind}--\ref{sec:magnetar-kind} specify
each kind in turn. Section~\ref{sec:verifier-interface} defines
the \texttt{P3QVerifier} Go interface that registers each kind
into the dispatch table. Section~\ref{sec:composition} ties the
three kinds back into LP-217 cert profile modes and shows the
rollup-contract composition pattern for the PQ-heavy posture.
